\chapter{Conclusion \& Future Work}
\hl{make the last point of the synthesis first point of te }
\section{Concluding Remarks}
\textit{Across all four contributions, what is the single most important thing that changes in how we should approach charging infrastructure planning at large geographic scales?}
The four analyses lead to clear insights to the geographic allocation of charging infrastructure for battery-electric vehicle fleets. These insights circle back to the three interconnected dimensions mentioned in Introduction (Section \dots, i-iii). For each of the interconnected dimensions, a key insight is raised, and one that is increasingly :

\begin{itemize}
    \item[(i)] \textbf{The identification of locations and required capacities for charging infrastructure:} A coherent implication of all contributions is the need to increasingly represent the operational patterns for charging infrastructure planning. Applied methodologies that go beyond the regional geographic scope are mostly focused on the capacity investment and do not typically represent the endogenous dynamics between capacity utilization, vehicle operation and optimal sizing. The co-optimization of charging infrastructure, charging activity, and BEV adoption in contributions 2 and 3 indicates that the operational patterns cannot be further considered in this \textit{stiff} manner and indicates the potential to increase the efficiency in infrastructure planning investments. \textit{from CAPEX to OPEX-centered} 
    
    \item[(ii)] \textbf{The influence of spatial configuration on charging demand and driver behavior}: Moreover, our findings emphasize the importance of defining expansion goals that are not solely based on the capacity or charging point per BEV, but complementary requirements that are based on maximum spacing within the road transport network. Currently, the AFIR defines a BEV count-based definition at the country level and distance-based for highways.
    
    \item[(iii)] \textbf{The integration of charging operations with the electricity system}: Electricity grid operators and regulators should leverage the potential for the spatial shifts. This inherently requires the consideration of the spatio-temporal movements of vehicles. Without this extended perspective, important 
    
    \item Overarchingly, infrastructure planning and policy design require cross-border collaboration in the charging infrastructure capacity planning, considering international and trans-regional transport flows. 
   
\end{itemize}


% \begin{itemize}
%     \item (core conclusion from paper3) When contextualizing this in previous works of system-level optimizations that go beyond charging infrastructure capacity planning and consider intangible costs based on the value of time parameter, the following is concluded: of \citet{LUH2024122412} and \citet{TATTINI2018265}; This highlights the importance of accounting for income-related differences in consumer behavior when planning charging infrastructure and developing optimization frameworks for passenger car electrification. The particular feature that needs to be accounted for is the income-dependent willingness to detour and spend time recharging. The results underscore the importance of planning public charging infrastructure with spatial equity considerations to ensure all neighborhoods have sufficient access. Moreover, our findings emphasize the importance of defining expansion goals that are not solely based on the capacity or charging point per BEV, but complementary requirements that are based on maximum spacing within the road transport network. Currently, the AFIR defines a BEV count-based definition at the country level and distance-based for highways \cite{transport2024afir}. 
    
%     \item (core conclusion paper 3) Consideration of income-class related detouring for charging
    
%     \item Formulation of spatial density-related policies: Moreover, our findings emphasize the importance of defining expansion goals that are not solely based on the capacity or charging point per BEV, but complementary requirements that are based on maximum spacing within the road transport network. Currently, the AFIR defines a BEV count-based definition at the country level and distance-based for highways \cite{transport2024afir}.

%     \item (paper 3) 
%     \item electricity grid opertors, und regulators: design von lokalen preissignalen zum leveragen von den spatial flexibility; cross-border cooperation 

% \end{itemize}
% \paragraph*{Consideration of operational patterns}

% \hl{Modal shift potential mehr reinbringen (HA)}
